\section*{Bab \ref{ch:b1:acetals-protection} --- Senyawa Karbonil: Adisi Nukleofilik, Asetal, dan Proteksi}

\begin{solution}{exo:b1:acetals-protection:1}
\ce{CH3CH(OH)OCH3}: \omterm{def:b1:acetals-protection:hemiacetal}{hemiasetal}. \ce{CH3CH(OCH3)2}: \omterm{def:b1:acetals-protection:hemiacetal}{asetal}.
\ce{CH3CH(OH)2}: hidrat. \ce{CH3OCH3}: eter. \ce{(CH3)2C(OCH2CH3)2}:
\omterm{def:b1:acetals-protection:hemiacetal}{asetal} (dari sebuah keton).
\end{solution}

\begin{solution}{exo:b1:acetals-protection:2}
\ce{CH3CH2CHO + 2CH3OH <=> CH3CH2CH(OCH3)2 + H2O}; propanon dengan
etana-1,2-diol menghasilkan 2,2-dimetil-1,3-dioksolana dan air,
\ce{(CH3)2CO + HOCH2CH2OH <=> C5H10O2 + H2O}.
\end{solution}

\begin{solution}{exo:b1:acetals-protection:3}
Natrium hidroksida, metilmagnesium bromida, dan natrium borohidrida
membiarkannya tidak berubah; asam sulfat encer berair menghidrolisisnya.
\end{solution}

\begin{solution}{exo:b1:acetals-protection:4}
OH pada karbon 4 beradisi pada karbon aldehida 1: cincin beranggota lima
yang terdiri atas empat karbon dan satu oksigen, dengan sebuah OH pada
karbon di sebelah oksigen cincin.
\end{solution}

\begin{solution}{exo:b1:acetals-protection:5}
Protonasi C=O; metanol beradisi pada karbon; lepasnya \ce{H+}: \omterm{def:b1:acetals-protection:hemiacetal}{hemiasetal};
protonasi OH-nya; lepasnya air menjadi \ce{CH3CH=O+CH3}; metanol kedua
beradisi; lepasnya \ce{H+}: \ce{CH3CH(OCH3)2}. Kedua proton yang diambil
dikembalikan: asam itu adalah \omterm{def:b1:elementary-steps:catalyst}{katalis}.
\end{solution}

\begin{solution}{exo:b1:acetals-protection:6}
Protonasi salah satu OCH3; lepasnya metanol menjadi \ce{(CH3)2C=O+CH3};
air beradisi; lepasnya \ce{H+}: \omterm{def:b1:acetals-protection:hemiacetal}{hemiasetal}; protonasi, lepasnya metanol,
lepasnya \ce{H+}: propanon. Air yang berlebih menjaga $Q$ di bawah $K$
untuk hidrolisis, yang berlangsung sempurna.
\end{solution}

\begin{solution}{exo:b1:acetals-protection:7}
Satu molekul air per \omterm{def:b1:acetals-protection:hemiacetal}{asetal}: $0.250 \times 18.0 = \qty{4.5}{g}$, sekitar
\qty{4.5}{mL}. Sesudah satu jam, $3.2/18.0 = \qty{0.178}{mol}$: konversi
\qty{71}{\percent}.
\end{solution}

\begin{solution}{exo:b1:acetals-protection:8}
Dengan diol, satu molekul menggantikan dua: reaksinya tidak mengurangi
jumlah molekul (satu aldehida dan satu diol menghasilkan satu \omterm{def:b1:acetals-protection:hemiacetal}{asetal} dan
satu air), dan OH kedua berada dekat dengan karbon yang bereaksi, sehingga
cincinnya mudah menutup. Keduanya membuat \omterm{def:b1:acetals-protection:hemiacetal}{asetal} siklik lebih
menguntungkan.
\end{solution}

\begin{solution}{exo:b1:acetals-protection:9}
Asam itu akan merusak \omterm{def:b1:organomagnesium:organometallic}{pereaksi Grignard} (perpindahan proton); jejak air
dan diol juga harus dihilangkan.
\end{solution}

\begin{solution}{exo:b1:acetals-protection:10}
1. Lindungi keton 4-klorobutan-2-on sebagai \omterm{def:b1:acetals-protection:hemiacetal}{asetal} siklik (etana-1,2-diol,
\omterm{def:b1:elementary-steps:catalyst}{katalis} asam, Dean--Stark). 2. Magnesium dalam eter kering: \omterm{def:b1:organomagnesium:organometallic}{pereaksi
Grignard} dari klorida yang terlindungi, yang tidak lagi dirusak oleh
karbonilnya sendiri. 3. Tambahkan propanon; pengolahan akhir dengan
amonium klorida berair: alkohol tersier. 4. Asam berair encer:
\omterm{def:b1:acetals-protection:protecting-group}{deproteksi} menjadi 5-hidroksi-5-metilheksan-2-on.
\end{solution}

\begin{solution}{exo:b1:acetals-protection:11}
$Q = [\text{asetal}][\ce{H2O}]/([\text{aldehida}][\ce{CH3OH}]^2)$: metanol
yang sangat berlebih, yang dikuadratkan dalam penyebut, menjaga $Q < K$
sampai aldehida yang tersisa tinggal sedikit. Pengeluaran air bekerja
pada pembilang; memakai metanol sebagai pelarut lebih sederhana, tetapi
memerlukan produk yang mudah dipisahkan.
\end{solution}

\begin{solution}{exo:b1:acetals-protection:12}
Hanya OH \omterm{def:b1:acetals-protection:hemiacetal}{hemiasetal} yang dapat pergi sebagai air dan menghasilkan kation
yang distabilkan oleh oksigen cincin di sebelahnya (\ce{C=O+}); gugus OH
yang lain harus pergi dari karbon biasa, yang menghasilkan kation yang
tidak terstabilkan. Metanol beradisi pada kation yang stabil itu: sebuah
metil \omterm{def:b1:acetals-protection:hemiacetal}{asetal}.
\end{solution}

\begin{solution}{pb:b1:acetals-protection:1}
\textbf{1.} \ce{BrMgCH2CH2CH2CHO}.
\textbf{2.} Karbonnya yang mirip \omterm{def:b1:electronic-effects:intermediates}{karbanion} akan menyerang aldehida molekul
lain (atau dirinya sendiri): \ce{RMgBr} beradisi pada $\mathrm{R'CHO}$, menghasilkan \omterm{def:b1:alcohol-activation:alkoxide}{alkoksida}.
\textbf{3.} Aldehidanya, terhadap \omterm{def:b1:organomagnesium:organometallic}{pereaksi Grignard}.
\textbf{4.} \omterm{def:b1:acetals-protection:hemiacetal}{Asetal} stabil terhadap \omterm{def:b1:organomagnesium:organometallic}{pereaksi Grignard} dan basa, dan dapat
dilepas oleh asam berair.
\textbf{5.} \omterm{def:b1:acetals-protection:hemiacetal}{Asetal} itu terbentuk lebih sempurna (satu diol untuk dua
molekul alkohol) dan mudah dihidrolisis.
\textbf{6.} $15.09/150.9 = \qty{0.1000}{mol}$.
\textbf{7.} $1.20 \times 0.1000 \times 62.0 = \qty{7.44}{g}$.
\textbf{8.} \ce{BrCH2CH2CH2CHO + HOCH2CH2OH <=> C6H11BrO2 + H2O}
(2-(3-bromopropil)-1,3-dioksolana).
\textbf{9.} \omterm{def:b1:elementary-steps:catalyst}{Katalis} itu mengaktifkan karbonil dan memungkinkan lepasnya
air pada tahap kedua; \omterm{def:b1:elementary-steps:catalyst}{katalis} itu terbentuk kembali.
\textbf{10.} Toluena dan air terdistilasi bersama, mengembun, dan jatuh ke
dalam penampung; air tenggelam dan tinggal di sana, toluena meluap kembali.
\textbf{11.} $Q$ mengandung $[\ce{H2O}]$ pada pembilangnya; pengeluaran air menjaga
$Q < K$, dan reaksinya berlanjut.
\textbf{12.} Air lebih rapat dan tidak bercampur dengan toluena: air
terkumpul di dasar; hanya lapisan atas yang mencapai lengan balik.
\textbf{13.} Ketika air berhenti terkumpul, pada volume yang diharapkan.
\textbf{14.} $0.1000 \times 194.9 = \qty{19.5}{g}$.
\textbf{15.} \ce{C6H11BrO2 + Mg -> C6H11BrMgO2} (eter kering).
\textbf{16.} Pereaksi itu beradisi pada \ce{(CH3)2CO}: \omterm{def:b1:alcohol-activation:alkoxide}{alkoksida}
\ce{(CH3)2C(O^-)CH2CH2CH2-} pada rantai yang terlindungi, dengan \ce{MgBr+}.
\textbf{17.} Asam akan melepas \omterm{def:b1:acetals-protection:protecting-group}{gugus pelindung} terlalu dini dan dapat
mendehidrasi alkohol tersier itu.
\textbf{18.} $0.0700 \times 174.0 = \qty{12.2}{g}$.
\textbf{19.} \omterm{def:b1:acetals-protection:hemiacetal}{Asetal} + \ce{H2O} $\to$ aldehida + etana-1,2-diol (\omterm{def:b1:elementary-steps:catalyst}{katalis}
asam).
\textbf{20.} Lepasnya air dari alkohol itu memerlukan \omterm{def:b1:predominant-reaction:strong-acid}{asam kuat} dan kalor;
asam berair yang lunak dan dingin menghidrolisis \omterm{def:b1:acetals-protection:hemiacetal}{asetal} jauh lebih cepat.
\textbf{21.} OH karbon 5 beradisi pada karbon aldehida 1: cincin
beranggota enam (lima karbon dan satu oksigen) dengan dua gugus metil pada
karbon di sebelah oksigen cincin dan sebuah OH pada karbon lain di
sebelahnya.
\textbf{22.} $0.0700 \times 0.90 = \qty{0.0630}{mol}$, \qty{8.19}{g}.
\textbf{23.} $0.0630/0.1000 = \qty{63}{\percent}$ (\omterm{def:b1:acetals-protection:protecting-group}{proteksi} dianggap
sempurna).
\textbf{24.} $0.1000 \times 18.0 = \qty{1.80}{g}$.
\textbf{25.} $1.80/0.997 =$ \textbf{\qty{1.8}{mL} air} dalam perangkap.
\end{solution}
